\bigskip\Needspace{13\baselineskip}
\pdfbookmark[1]{Tables for the Perturbations of Ceres. March 1803}{ceresstoerungen180303}
\begin{center}
TABLES FOR THE PERTURBATIONS OF CERES.\par
\medskip\hrule\medskip
Monatliche Correspondenz. Vol. VII, pp.\ 259—262. March 1803.\par
\medskip\hrule\bigskip
GAUSS TO VON ZACH.
\end{center}\medskip
\begin{flushright}Braunschweig, 26 December 1802.\end{flushright}

Some time ago I wrote to you that I had been thinking about shortening the tables for the perturbations of \textit{Ceres} and intended to send you a sample later. (Mon. Corr. vol. VI, p.\ 498) [p.\ 230 of this volume] I now have the honor of explaining the matter more fully and sending you a table \textit{in extenso} in place of a sample.

Even at this point, with my calculations still limited to the perturbations due to Jupiter and to the first powers of the eccentricities, there are so many equations that it would be extremely laborious to calculate their numerical values for a large number of positions from the formulas themselves without tables. Even with tables arranged in the usual way, however, the work really is laborious. For \textit{all} the equations printed in the November 1802 issue of the M.\ C. [p.\ 228 of this volume], \textit{forty} tables would be needed, and nearly \textit{thirty} even if the smaller equations that amount individually to less than $2''$ were omitted. Forming so many arguments, looking them up in so many tables and making so many additions still takes a great deal of patience if it has to be done often.
I have therefore tried to shorten these tables in a more convenient way. The first \textit{eight} equations \textit{for longitude} (the first coefficient on p.\ 495 [228], incidentally, has to change from $230''00$ to $231''00$) naturally go into \textit{one} table, whose argument is $\CeresSymbol-\JupiterSymbol$. The other fourteen, however, would require as many tables. They could be put into a double-entry table, but it would be so large, tedious to calculate and cumbersome to use as to be out of the question. I have therefore chosen the following procedure:
\SourceLeaf{241}%

If I write $\CeresSymbol-\JupiterSymbol$ as $D$ for brevity, with a little attention one can see that each of these fourteen equations can be put into the form
\[
a\sin(iD-\CeresSymbol-A)=a\sin(iD-A)\cos\CeresSymbol-a\cos(iD-A)\sin\CeresSymbol
\]
where $i$ represents positive and negative integers. The sum of all the equations will therefore be expressed by $P\cos\CeresSymbol-Q\sin\CeresSymbol$, where $P$ and $Q$ are functions of $D$. If one sets
\[
\frac PQ=\operatorname{tang}\varphi,\quad
\frac P{\sin\varphi}=\frac Q{\cos\varphi}=R,
\]
then $R$ and $\varphi$ also depend only on $D$, and the sum of all fourteen equations
\[
=R\sin(\varphi-\CeresSymbol)
\]
Thus in a table with argument $D$, for every value of $D$ one can enter the corresponding values of $R$ and $\varphi$ and immediately obtain the combined value of all those equations by a very easy calculation. I have, however, decided against tabulating $\varphi$ and entered $\varphi-2D$ instead, because this angle varies less than $\varphi$ and only oscillates periodically within certain limits as $D$ takes all values from $0$ to $360^0$, whereas $\varphi$ would make two revolutions in that time. This $\varphi-2D$ is denoted by $B$ in the third table, just as the $R$ above is denoted by $A$. In this way, the combined value of all the equations for longitude that depend on the first power of the eccentricity is
\[
A\sin(B+2D-\CeresSymbol)=A\sin(B+\CeresSymbol-2\JupiterSymbol)
\]

Similarly, the equations for the radius vector that are independent of the eccentricity are included in Table II. Those that depend on the first powers of the eccentricities are put into the form $A\sin(B+2\CeresSymbol-3\JupiterSymbol)$ in Table IV. Finally, Table V puts into that same form all the equations for latitude listed in the November issue [p.\ 229 of this volume].

In this way we have only five tables. This does call for some trigonometric calculation, but I at least find it far more convenient than having to look up each of the equations combined here in a separate table. Some further changes could also be made to advantage, for example by replacing $A$ directly with $\log A$, \textit{or by some other means}. But I shall leave this for the future, since these tables are only meant to serve until future observations have led to further corrections of the elements and the perturbations have accordingly been calculated more accurately and more completely. For now I think these tables should serve me well, especially when I calculate next year’s observations.

\Needspace{7\baselineskip}
\begin{center}\textit{Remarks.}\end{center}
1.\quad $\CeresSymbol$ and $\JupiterSymbol$ denote the mean heliocentric longitudes of Ceres and Jupiter.

2.\quad The perturbations of the radius vector are expressed in millionths of the radius of the Earth’s orbit.

3.\quad In the perturbations of latitude, and in the second parts of the perturbations of longitude and of the radius vector, $A$ is to be treated as positive, and due attention must be paid to the sign of the sine of $B+\CeresSymbol-2\JupiterSymbol$ and of $B+2\CeresSymbol-3\JupiterSymbol$.

4.\quad For latitude, take $B$ from the first column for the first six signs of the argument, or when descending. Take it from the third column in the last six signs, or when ascending. Both use the same $A$. The arrangement of the table makes this clear.
